% D2-1Agent
% Auto-generated from markdown source
% origin: ai

\documentclass[12pt, a4paper]{article}
\usepackage[utf8]{inputenc}
\usepackage[T1]{fontenc}
\usepackage{amsmath, amssymb, amsfonts, amsthm}
\usepackage{textcomp}
\usepackage{geometry}
\usepackage{hyperref}
\usepackage{booktabs}
\usepackage{graphicx}
\usepackage{fancyhdr}
\usepackage{tcolorbox}
\usepackage{enumitem}

\geometry{margin=1in}
\pagestyle{fancy}
\fancyhf{}
\rhead{D2-1Agent}
\lhead{\thepage}

\newtheorem{definition}{Definition}[section]
\newtheorem{theorem}{Theorem}[section]
\newtheorem{corollary}{Corollary}[section]
\newtheorem{proposition}{Proposition}[section]
\newtheorem{lemma}{Lemma}[section]
\newtheorem{conjecture}{Conjecture}[section]

\newcommand{\calI}{\mathcal{I}}
\newcommand{\calP}{\mathcal{P}}
\newcommand{\calE}{\mathcal{E}}
\newcommand{\calS}{\mathcal{S}}
\newcommand{\calH}{\mathcal{H}}
\newcommand{\calO}{\mathcal{O}}
\newcommand{\calF}{\mathcal{F}}
\newcommand{\calU}{\mathcal{U}}
\newcommand{\calV}{\mathcal{V}}
\newcommand{\calA}{\mathcal{A}}
\newcommand{\calG}{\mathcal{G}}
\newcommand{\calB}{\mathcal{B}}
\newcommand{\calC}{\mathcal{C}}
\newcommand{\calL}{\mathcal{L}}
\newcommand{\calM}{\mathcal{M}}
\newcommand{\calX}{\mathcal{X}}
\newcommand{\calY}{\mathcal{Y}}
\newcommand{\calW}{\mathcal{W}}
\newcommand{\calT}{\mathcal{T}}
\newcommand{\calZ}{\mathcal{Z}}
\newcommand{\calR}{\mathcal{R}}
\newcommand{\calQ}{\mathcal{Q}}
\newcommand{\calD}{\mathcal{D}}
\newcommand{\calN}{\mathcal{N}}
\newcommand{\calK}{\mathcal{K}}

\newcommand{\Tr}{\text{Tr}}
\newcommand{\E}{\mathbb{E}}
\newcommand{\diff}{\mathrm{d}}
\newcommand{\KL}{D_{\mathrm{KL}}}
\newcommand{\argmax}{\arg\max}
\newcommand{\argmin}{\arg\min}
\newcommand{\hbarF}{\hbar}

\newtcolorbox{thmbox}[1]{
  colback=blue!5!white, colframe=blue!75!black, fonttitle=\bfseries,
  title={#1}, breakable
}
\newtcolorbox{defbox}[1]{
  colback=green!5!white, colframe=green!75!black, fonttitle=\bfseries,
  title={#1}, breakable
}

\begin{document}

\title{D2-1Agent}
\author{RALPH Loop}
\date{}
\maketitle

\begin{abstract}
This paper presents D2-1Agent, formalizing the underlying principles using information-theoretic methods.

\textbf{Keywords}: D2-1Agent, information theory, formalization.
\end{abstract}

\tableofcontents
\section{D2-1Agent}




\#\#



\#\#

\textbf{Agent } \(A\)  Agent \(\mathcal{T}\) \(L\)****

\[I_{plan} = I(\text{plan}; \mathcal{T}, y_{target})\]

**** \(|A| < I_{plan}^* + \delta_{margin}\) \(\delta_{margin} \approx 10\%\) Agent  \((L, |\mathcal{T}|)\)


\#\#

\subsubsection{ 1}


\textbf{ vs }
\begin{itemize}
\item IB\(X \to Z \to Y\)
\item \textbf{ IB}\(\mathcal{T} \to \text{plan} \to \text{action\_sequence}\)
\end{itemize}


\textbf{ \(I_{plan}\) }
1. ****
2. ****" A  B "
3. ****

****
\[I_{plan} \geq I(\mathcal{T}; y_{target}) - I_{irrelevant}(\mathcal{T}; y_{target})\]


\subsubsection{ 2}


****
\[|A|_{crit} = I_{plan}^* + \delta_{margin}\]

\(\delta_{margin}\)  \(d\)
\[\delta_{margin} \geq d \cdot \log |\mathcal{T}|\]

****
\[\text{}(A) = \mathbb{I}\left[|A| > I_{plan}^* + \delta_{margin}\right]\]


\subsubsection{ 3}


\textbf{ \(|A|_c\)}\textbf{** \(|A|\) }**""""

****
\[\lim_{|A| \to |A|_c} \frac{d}{d|A|} P_{\text{planning}}(|A|) = \infty\]

\textbf{ grokking }
\begin{itemize}
\item Grokking
\end{itemize}

-
\begin{itemize}
\item ****\(\Psi(x) = \tanh(k(x - x_c))\)
\end{itemize}



\subsubsection{ 4 D9-III}


\textbf{D9-III  + }
\begin{itemize}
\item D9-III  Agent ****
\item D2-1  Agent ****
\end{itemize}


****
\[\text{} = \mathbb{I}[|A| > |A|_c] \land \mathbb{I}[I_{plan} > \mathbb{I}_{\text{cov}}(\mathcal{T})]\]

****D2-1 D9-III


\#\#

\subsubsection{ 1}

```python
\section{1.}

\section{2.  P\_planning vs |A|}

\section{3.  |A|\_c}

```

\subsubsection{ 2}

```python
\section{1.  L}

\section{2.}

\section{3.  L  |A|\_c}

```

\subsubsection{ 3}

```python
\section{1.  I\_plan}

\section{2.  I\_plan}

```


\#\#

1. ****
2. \textbf{""""}
3. \textbf{ grokking }


\#\#

\begin{itemize}
\item [x]  1-4
\item [ ]  1
\item [ ]  2
\item [ ]  3
\end{itemize}



* D2-1  D9-I *
*2026-05-04 | r01 | *

\begin{thebibliography}{99}
\bibitem{ref1} Shannon, C.E. (1948). A Mathematical Theory of Communication.
\bibitem{ref2} Cover, T.M. (2006). Elements of Information Theory.
\end{thebibliography}

\end{document}